% !TEX root = tarasov-qasp-zh.tex
% 章首扉图与格言（sec-09，分子物理学）在英文源 ch-17.tex 尾部（有效代码），已归 zh/ch-17-zh.tex；
% 本章首尾均无有效 sec 块（ch-19.tex 尾部那组 sec 块为 % 注释的死代码副本，未搬入）。

\chapter[你怎样解释水的热膨胀的特殊性？]{你怎样解释水的热膨胀的特殊性？\label{ch-19}}

\begin{dialogue}

\tchr 水的热膨胀有哪些特殊性？

\tchr 水从 \SI{0}{\degreeCelsius} 加热到 \SI{4}{\degreeCelsius} 时，密度增大。只有当温度升到 \SI{4}{\degreeCelsius} 以上时，它才开始膨胀。

\tchr 这一点你怎样解释？

\stda 我不知道。

\pnote{底本本章开篇连续三个回合均标“教师”：教师发问后随即自陈事实，再问“怎样解释”。对照\chap{ch-18} 开篇的问答节奏（教师提问$\to$学生甲复述事实$\to$教师追问$\to$学生甲答不出），第二个回合的说话者疑为学生甲，疑系重排时漏标说话人。本章参照 PDF 与底本同源，无从核实；中文版照底本译出，待考 1973 年原书。本章另有拼写级排印小误（如 ``on the her hand'' 应为 ``on the other hand''），径改不注。}

\begin{figure}[!h]%
\centering
\includegraphics[width=0.4\textwidth]{figs/ch-19/fig-072a}
\caption[解释水的热膨胀的特殊性。]{解释水的热膨胀的特殊性。\label{fig-72}}
\end{figure}

\tchr 水的这一独特性质与它的原子结构有关。水分子之间只能以一种方式相互作用：每个水分子只能再连上四个邻近的分子，这些分子的中心恰好构成一个四面体（\figr{fig-72}）。由此形成一种疏松的花边状结构，它表明水具有准晶性。当然，谈到水的结构，正如谈到任何别的液体一样，只能在短程的层面上来谈（见\chap{ch-18}）。随着与所选分子距离的增大，这种有序性会因分子间键的弯曲与断裂而逐渐畸变。温度升高时，分子间的键断裂得愈发频繁，带有未占用键的分子越来越多，它们填入四面体结构的空位，于是准晶性的程度随之降低。
\pnote{原书对这些“键”的实体留而未说：它们是氢键（hydrogen bond）。水分子中的氧原子有两对孤对电子，分子又带两个氢原子，因此每个水分子至多形成四个氢键，其方向恰好排成四面体，即\figr{fig-72} 所示的图像；冰正是这样开放的四面体骨架，密度只有约 $0.92~\si{\gram\per\cubic\centi\metre}$，反而小于液态水；液态水中的平均氢键数则不足四个。氢键在人教版高中化学选择性必修2《物质结构与性质》讲分子间作用力时介绍，“冰浮于水面”是其经典例证。“准晶性”一名的来历见\chap{ch-18} 评注；下文“两种竞争效应”的定性解释至今仍是标准说法。}

上面所说的水这种准晶物质的花边状结构，令人信服地解释了水的物理性质的反常，特别是其热膨胀的特殊性。一方面，温度升高使分子内的振动加剧，原子间与分子间的平均距离随之增大，也就是说，分子似乎略微“发胀”。另一方面，温度升高又会破坏水的花边状结构，这自然使分子本身堆积得更加紧密。第一种效应（振动效应）应使水的密度减小——这正是引起固体热膨胀的那种常见效应。第二种效应，即结构破坏效应，则相反，应使水受热时密度增大。把水加热到 \SI{4}{\degreeCelsius} 时，结构效应占上风，水的密度因而增大；继续加热，振动效应开始占上风，水的密度便随之减小。
\pnote{常压下水密度最大处实际在约 \SI{3.98}{\degreeCelsius}，本书与教材一律取整为 \SI{4}{\degreeCelsius}。从 \SI{0}{\degreeCelsius} 到 \SI{4}{\degreeCelsius} 的“热缩冷胀”即初中物理所讲水的反常膨胀（人教版八年级上册 6.4《密度与社会生活》）：冬夜湖面先结冰，而湖底水温保持在约 \SI{4}{\degreeCelsius}，水生生物得以越冬——密度最大的水恰好沉在湖底。}

\end{dialogue}


\cleardoublepage

% NEWTERMS: peculiarity | 特殊性 | 章题用词；anomaly（of physical properties）译“（物理性质的）反常”，中文教学习称“水的反常膨胀”
% NEWTERMS: thermal expansion of water | 水的热膨胀 | 承 ch-17 章尾格言首现（彼处未括注），本章为正文首现，不重复括注
% NEWTERMS: tetrahedron | 四面体 | 普通几何词，不括注；指四个邻近分子的中心所构成
% NEWTERMS: friable lace-like structure | 疏松的花边状结构 | 原书描述性用语，照译
% NEWTERMS: quasi-crystallinity | 准晶性 | 承 ch-18 定名（含英文括注与勿混淆“准晶”之评注），本章沿用不重复括注
% NEWTERMS: short-range (level) | 短程（层面） | 承 ch-18“短程有序”
% NEWTERMS: intermolecular bonds | 分子间键 | 实体即氢键，见评注
% NEWTERMS: unoccupied bonds / vacancies | 未占用键／空位 | 原书描述性用语照译
% NEWTERMS: intramolecular vibrations | 分子内的振动 | 承 ch-18“分子内原子的振动”
% NEWTERMS: vibrational effect / structural (breakup) effect | 振动效应／结构（破坏）效应 | 原书自铸的一对对立短语
% NEWTERMS: hydrogen bond | 氢键 | 仅评注层引入；人教版高中化学选择性必修2《物质结构与性质》
